\documentclass[kant]{rhelder3}

\title{The Paralogisms of Practical Reason Define the Phenomena}
\author{%
  \initials{R.W.} Helder%
  \thanks{%
    I want to thank God,
    because God gave me this Dundie,
    and I feel God in this Chili's tonight.%
  }%
  \\\acro{UC} Berkeley\and
  Immanuel Kant\\University of Königsberg%
}
\date{August 3, \digits{2026}}

\begin{document}
\maketitle

\begin{abstract}
  \autokant
\end{abstract}

\kantsection
\autokant[2]

\kantsubsection
\autokant[2]

\kantsubsubsection
\autokant[2]

\kantparagraph
\autokant

\kantsubparagraph
\autokant*\unskip\footnote{\autokant[1][1]}

\section{Nested Sections}
\unskip\vspace{12pt}
\subsection{Nested Subections}
\subsubsection{Nested Subsubsections}

\autokant*\unskip\footnote{\autokant[0][1]}

\autokant

\section{Block Quotes}

\autokant*
\begin{quote}
  \autokant*[1][2]
\end{quote}
\autokant

\section{Lists}

\subsection{Enumerate}

\autokant*
\begin{enumerate}
  \item \autokant*[1][1]
  \item \autokant*[0][1]
  \item \autokant*[0][1]
\end{enumerate}
\autokant

\subsection{Itemize}

\autokant
\begin{itemize}
  \item \autokant*[1][1]
  \item \autokant*[0][1]
  \item \autokant*[0][1]
\end{itemize}
\autokant

\subsection{Description}

\autokant
\begin{description}
  \item[The Ideal] \autokant*[1][1]
  \item[Practical Reason] \autokant*[0][1]
  \item[The Categories] \autokant*[0][1]
\end{description}
\autokant

\section{Nested Lists}

\subsection{Enumerate}

\autokant
\autokant*
\begin{enumerate}
  \RaggedRight
  \item \autokant*[1][1]
    \begin{enumerate}
      \item \autokant*[0][1]
        \begin{enumerate}
          \item \autokant*[0][1]
            \begin{enumerate}
              \item \autokant*[0][1]
              \item \autokant*[1][1]
            \end{enumerate}
          \item \autokant*[0][1]
        \end{enumerate}
      \item \autokant*[0][1]
    \end{enumerate}
  \item \autokant*[0][1]\unskip\footnote{\justifying\autokant[0][1]}
\end{enumerate}
\autokant

\subsection{Itemize}

\autokant
\autokant*
\begin{itemize}
  \RaggedRight
  \item \autokant*[1][1]
    \begin{itemize}
      \item \autokant*[0][1]
        \begin{itemize}
          \item \autokant*[0][1]
            \begin{itemize}
              \item \autokant*[0][1]
              \item \autokant*[1][1]
            \end{itemize}
          \item \autokant*[0][1]
        \end{itemize}
      \item \autokant*[0][1]
    \end{itemize}
  \item \autokant*[0][1]
\end{itemize}
\autokant


\section{Secondary Fonts}

\subsection{Sans Serif}

\autokant*[1][1]
\textsf{\autokant*[0][1]}
\autokant*[0][1]

\textsf{\autokant[1][3]}

\autokant[1][2]

\subsection{Typewriter}

As is evident upon close examination,
\texttt{let us suppose that},
in accordance with the principles of \texttt{time},
our a priori concepts are the clue to the discovery of philosophy.
\texttt{Metaphysics exists in the objects in space and time.}
\autokant*[1][1]

\subsection{Math}

\autokant*[1][1]
\[ \sum_{0\leq k<n} k = \frac{n(n-1)}{2} \]
\autokant*[0][1]
\[ \sin \frac{\alpha}{2} = \pm \sqrt{\frac{1-\cos\alpha}{2}} \]
\autokant*[0][1]

\autokant*[1][1]
\[ \sin \frac{\alpha}{2} = \pm \sqrt{\frac{1-\cos\alpha}{2}} \]
\autokant*[1][2]
\(\sum_{i=1}^n\), \(\int_0^\infty\), \(\lim_{n \to 0}\)
\autokant*[0][2]

\section{Quotations and Ellipses}

\textquote[][.]{The hair is gray, the hay is folly}

\textquote{The sweets are cherry, the lads are merry}.

\textquote[A Loon]{Bye and by and forth to go},
\textquote[Loon Two][!]{Nary a win is just for show}

One by one \textellipsis{} and two by two.\thinspace\textellipsis{}
Trusty bulworks \textellipsis, and crackers too!

\enquote{A \textelp*{lie} for an eye is a \textelp{stye} on the sly \textelp{}}

\subsection{Latin}

The \texttitle{Aeneid} begins with
\foreigntextquote{latin}{Arma virmumque cano},
which for some reason confused me as a high schooler;
after all, why sing of his arms, and not his legs?
Perhaps Virgil was not being so serious:
\foreigntextquote{latin}{\foreigntextquote{latin}{Arma} virmumque cano}.

\subsection{Greek}

Aristotle perplexingly says
\foreigntextquote{ancientgreek}{%
  ἐκ τῶν καθ' ἕκαστα γὰρ τὰ καθόλου·
  τούτων οὖν ἔχειν δεῖ αἴσθησιν, αὕτη δ' ἐστὶ νοῦς%
},
which actually means that \alialingua{nous}
is perception of universals.
Perhaps Aristotle did not mean it seriously:
\foreigntextquote{ancientgreek}[][.]{%
  ἐκ τῶν καθ' ἕκαστα γὰρ \foreigntextquote{ancientgreek}{τὰ καθόλου}·
  τούτων οὖν ἔχειν δεῖ αἴσθησιν, αὕτη δ' ἐστὶ νοῦς%
}

\end{document}
